
\subsubsection{3.1 Model Set}

The analysis is conducted over N = 30 LLMs spanning 9 families: LLaMA (9 models), Mistral (6), Qwen (6), Gemma (2), Falcon (2), SOLAR (2), MPT (1), StableLM (1), and Phi (1). The set covers 3 parameter scales (7B, 13B, 70B+) and 2 training regimes (pretrained-only, instruction-tuned/RLHF).

\subsubsection{3.2 Adversarial Fragility: AdvGLUE Accuracy Drop}

Adversarial fragility is operationalized as AdvGLUE accuracy drop — the difference between benign and adversarial accuracy across 14 attack methods on GLUE-derived NLU tasks. Direct AdvGLUE measurements from TrustLLM ICML 2024 Table 2 were available for 11 of 30 models. For the remaining 22 models, AdvGLUE drop was estimated via OLS regression trained on the 11 anchor measurements, with each estimated value flagged in the dataset. The TrustLLM HuggingFace dataset required user agreement (HTTP 403 error) and was inaccessible; published paper values were used as anchors. This imputation procedure is a material limitation (see Section 6.2, L1): 73\% of AdvGLUE values are estimated rather than directly measured, and OLS imputation may compress variance and introduce correlation structures in RI.

AdvGLUE drop values in the assembled dataset range across the 30 models with SD = 0.1212 (95\% CI: [0.093, 0.138]).

\subsubsection{3.3 Capability Index: PC1}

PCA is applied to six Open LLM Leaderboard v2 benchmark scores — BBH, ARC-Challenge, MMLU-Pro, MATH, GPQA, and MuSR — to produce a single capability composite (PC1). PC1 explains 68.5\% of benchmark variance across the 30 models. This falls marginally below the pre-registered 70\% threshold, due to the harder and more diverse tasks in Leaderboard v2. The shortfall is recorded as a sensitivity note (L3, Section 6.2).

ARC-Challenge contributes to PC1 construction (Section 3.3) and is also the source benchmark for ECE measurement (Section 3.5). OLS residualization removes the linear PC1–ECE relationship by construction, but benchmark-specific features of arc_challenge may still influence the RI–ECE residual correlation. This potential circularity is noted as limitation L7.

\subsubsection{3.4 Residual Instability (RI)}

RI is defined as the OLS residual from the regression:

```
AdvGLUE_drop ~ PC1 + mean_confidence
```

where \texttt{mean_confidence} is the mean of per-sample maximum-choice softmax probabilities from arc_challenge (range across models: 0.789–0.958). By OLS construction, RI is orthogonal to PC1. Empirical VIF = 1.000 confirms orthogonality. The RI construct passes two pre-registered gate conditions: SD(AdvGLUE_drop) = 0.1212 > 0.05 (PASS) and R²_residualization = 0.529 < 0.80 (PASS). The R² value indicates that capability (PC1) and mean confidence together explain 52.9\% of adversarial drop variance, leaving substantial model-specific residual.

\subsubsection{3.5 Calibration Measurement: ECE}

Expected Calibration Error is computed from per-sample softmax probabilities extracted from arc_challenge log-likelihoods (Open LLM Leaderboard v2, n = 1,172 samples per model), using 10-bin equal-width binning. ECE values across the 30 models range from 0.175 (meta-llama/Meta-Llama-3-70B) to 0.472 (stabilityai/stablelm-zephyr-3b).

\subsubsection{3.6 Statistical Analysis}

The primary analysis is the Spearman partial correlation ρ(RI, ECE | PC1, mean_confidence), computed via \texttt{pingouin.partial_corr()} \citep{Vallat2018}. Holm-Bonferroni correction is applied across 4 pre-registered predictions. Within-family analyses are conducted for the three largest families: LLaMA (n = 9), Mistral (n = 6), and Qwen (n = 6). Bootstrap confidence intervals use 10,000 resamples with seed 42. VIF is computed for all covariates; Cook's distance is computed for outlier identification. All computations are CPU-based (Python 3.10; scikit-learn 1.4.2; pingouin 0.6.1; scipy 1.13.0; statsmodels 0.14.2; uncertainty-calibration 0.1.4). The pipeline is validated by 41/41 unit tests.

